\chapter{UVOD}
\label{ch:uvod}

TODO: piše se PRETPOSLJEDNJE (prije zaključka), kad su rezultati poznati.

\section{Motivacija}
\label{sec:motivacija}

Industrijska robotika od svog začetka počiva na pretpostavci strukturirane okoline. Robot je
ograđen zaštitnom ogradom, radni komadi nalaze se na predodređenom mjestu, a zadatak
se svodi na ponavljanje prethodno naučene i isprogramirane radnje. U takvim uvjetima pozicijsko
upravljanje visoke krutosti daje izvrsne rezultate, robot vrlo ponovljivo i bez većih odstupanja
slijedi zadanu putanju bez obzira na smetnje.

Ta se pretpostavka ruši čim robot izađe iz kaveza. Industrijska okruženja puna su ljudi koji istovremeno
obavljaju svoje poslove, servise, logističke zadatke i okolnosti koje se mijenjaju između dva izvođenja istog zadatka.
Odgovor na takve probleme su mobilni manipulatori - sustavi koji integriraju svestranost mobilne platforme
s artikuliranošću robotske ruke i time uklanjaju ograničenje dosega koje veže nepokretni manipulator za jedno radno mjesto.

Najčešća primjena industrijskih robota je za tzv. \textit{pick-and-place} radnje, ali postoji
i veliki broj operacija za koje je potrebna uspostava i održavanje fizičkog kontakta, primjerice
otvaranje vrata i ladica, umetanje predmeta, posluživanje stroja ili pomicanje objekata.
Postojanje kontakta uvelike mijenja problem upravljanja. Dok se robot giba u slobodnom prostoru,
greška položaja je relativno mala i nema brige da će robot oštetiti okolinu.
U kontaktu s krutom okolinom ista greška položaja proizvodi silu koja raste s krutošću regulatora
i može oštetiti predmet, okolinu ili samog robota. Zbog toga upravljanje u kontaktu ne može biti
čisto pozicijsko, već mora regulirati odnos između položaja i sile \cite{raibert1981hybrid}.

Otvaranje vrata reprezentativan je predstavnik te klase zadataka i zbog toga se u literaturi
često koristi kao referentni problem \cite{ott2005door}. Vrata su artikulirani objekt čije je gibanje ograničeno
na jedan stupanj slobode, ali robot taj stupanj slobode ne poznaje unaprijed.
Drugim riječima, robot ne zna gdje je os šarke, u kojem smjeru se krilo otvara ili koliki je otpor u zglobu.
Sve to mora zaključiti iz vlastitih mjerenja tijekom interakcije \cite{sturm2011probabilistic}.
Uz to, hod krila redovito premašuje doseg ruke pri nepomičnoj
bazi, pa zadatak traži i koordinaciju platforme i manipulatora. Zadatak time istovremeno
obuhvaća percepciju, procjenu svojstava okoline i upravljanje u kontaktu, što ga čini prikladnim
ispitnim slučajem za cjelovit sustav.

Postoje dva prevladavajuća pristupa takvim zadacima. Klasični pristup temeljen na
modelu oslanja se na eksplicitan opis zadatka: procijeni se geometrija ograničenja, iz nje se
generira referentna putanja, a impedancijski ili hibridni silno-pozicijski regulator izvodi
gibanje uz propisanu popustljivost \cite{khatib1987operational}. Prednost je razumljivost i
predvidljivost, svaki je korak moguće analizirati i obrazložiti. Slabost je osjetljivost na
kvalitetu procjene: pogreška u procijenjenoj geometriji ograničenja izravno se pretvara u
kontaktnu silu. Pristup temeljen na učenju ne zahtijeva eksplicitan model, već strategiju
upravljanja stječe kroz interakciju sa simuliranom okolinom. Novije formulacije akcijskim
prostorom politike čine samu krutost regulatora, pa politika uči ne samo kamo se gibati nego i
koliko biti popustljiva u pojedinom trenutku \cite{martinmartin2019variable}. Cijena je gubitak
uvida u razloge odluke i ovisnost o vjernosti simulacije.

Usporedba tih dvaju pristupa ima smisla samo ako se provodi na istom zadatku, istom robotu i istoj fizici okoline.
Ovaj rad zato oba pristupa razvija za isti sustav (mobilni manipulator KUKA KMR iiwa u simulacijskom okruženju),
i uspoređuje ih na istim mjerilima uspješnosti.

\section{Struktura rada}
\label{sec:struktura}

Poglavlje \ref{ch:pregled} daje pregled dosadašnjih istraživanja, od klasičnih zakona
upravljanja u kontaktu preko procjene kinematičkih modela artikuliranih objekata do novijih
pristupa temeljenih na podržanom učenju. Poglavlje \ref{ch:sustav} opisuje robotski sustav i
simulacijsko okruženje: mobilni manipulator KMR iiwa, prihvatnicu izrađenu za ovaj zadatak,
modele vrata te programsku arhitekturu koja povezuje simulator i upravljački kod.

Poglavlja \ref{ch:percepcija} i \ref{ch:procjena} obrađuju ulazne module sustava. Prvo opisuje
određivanje položaja i orijentacije kvake iz slike, a drugo procjenu sile i momenta na vrhu
alata te postupak kojim sustav iz izmjerenog odziva zaključuje o ograničenju gibanja i krutosti
okoline.

Poglavlja \ref{ch:klasicni} i \ref{ch:rl} predstavljaju dva uspoređivana pristupa upravljanju.
Prvo opisuje rješenje temeljeno na modelu, organizirano kao slijed faza s eksplicitnom procjenom
geometrije ograničenja. Drugo opisuje formulaciju zadatka kao Markovljevog procesa odlučivanja,
prostor stanja i akcija s promjenjivom impedancijom, oblikovanje funkcije nagrade i postupak
učenja.

Poglavlje \ref{ch:rezultati} donosi eksperimentalne rezultate, ablacijske studije i usporedbu
dvaju pristupa na istoj fizici okoline, uz raspravu o ograničenjima provedene usporedbe.

Poglavlje \ref{ch:zakljucak} sažima postignuto i navodi smjernice za daljnji rad.